\subsection{Is one approach better than the other?}
\label{sec:better}
A natural question at this point from the language designer's perspective is: which approach to
linearity should I use in my language? We do not try to offer an authoritative answer, but instead give
some observations and points to consider.

Clearly, graded-base style approaches afford a simpler system for the end-user with less syntactic overhead. The
pervasive grading obviates the need for most promotions and \textsf{let}-unboxings.
But while having explicit \textsf{push} syntax is an additional burden to the user of a
linear-base system, in practice this is not necessary. For example, Granule handles this implicitly
via pattern matching~\cite{DBLP:journals/pacmpl/OrchardLE19, hughes2021linear}.

On the flip-side, it may be desirable to disallow \textsf{push} for reasoning about certain
domains, where a coeffect should not distribute over the components of a pair. An example for this
might be for reasoning about interdependent threads that need to run in parallel. Given what we
have shown in this paper, linear base systems would seem better equipped to offer this more restrictive
modality, however \citet{hughes2021linear} have shown that it is possible to recover the non-distributive behaviour of the $\Box$-modality with respect to products in a graded-base setting.
They argue that any useful language implementation would reasonably have to offer both flavours of
the modality.

What about the burden for the language designer? Empirically, the authors have found that proofs for
graded-base systems are significantly shorter, due, in part, to the homogeneous treatment of
assumptions, i.e. the lack of split between linear and non-linear free variables. Especially
when retrofitting linear or coeffect types to an existing language, only Graded Base offers a
backwards-compatible way and in such settings this is the only practical option (see, e.g.,
the addition of graded types to Haskell for linearity~\cite{DBLP:journals/pacmpl/BernardyBNJS18}).

\subsection{Graded Type Systems in Practice}
\label{sec:wild}
The \haskin{LinearTypes} language extension available since GHC 9.0.1~\cite{GHCLinearTypes}
enables Linear Haskell in the context of an
industrial-strength programming language. As discussed in this paper, Linear Haskell is a graded-base system, which is currently semiring-monomorphic: in the theory
of \citet{DBLP:journals/pacmpl/BernardyBNJS18} the \textit{none-one-tons} semiring is used,
$\{0, 1, \omega\}$. However, as of the time of writing,
the implementation in GHC offers no $0$, and
thus the grades range over $\{1, \omega\}$, written in Haskell as \haskin{One} and \haskin{Many}, respectively.
Consider the following Linear Haskell program:
%
\begin{haskell}
{-# LANGUAGE LinearTypes #-}
{-# LANGUAGE DataKinds #-}
import GHC.Types
const :: a %One -> b %Many -> a
const x y = x
\end{haskell}
%
By default arrows are \haskin{Many} graded, thus the above can also be written
as:
%
\begin{haskell}[firstnumber=last]
const' :: a %One -> b -> a
const' x y = x
\end{haskell}
%
Notably, Linear Haskell supports \textit{multiplicity polymorphism}, however the multiplicities
are inhabited by a single semiring.
Future work would be a \haskin{GradedTypes} extension to generalise Haskell to be polymorphic over
semirings, a feature supported by Granule, where \citet{DBLP:journals/pacmpl/OrchardLE19} use an
SMT-solver to discharge constraints over the grades. It is not yet clear how this could be
practically integrated with GHC, perhaps as a compiler plugin akin to Liquid Haskell \cite{Vazou2020LiquidHaskell}.

Idris 2 also implements a graded-base system, based on Quantitative
Type Theory (due initially to \citet{McBride2016} and later adapted
slightly by \citet{quantitative-type-theory}), also providing
grades in the \textit{none-one-tons} semiring~\cite{DBLP:conf/ecoop/Brady21,DBLP:journals/darts/Brady21}. For example, \haskin{const} can be rendered in Idris 2:
%
\begin{haskell}
const : {0 a : Type} -> {0 b : Type} -> a -> (0 x : b) -> a
const x y = x
\end{haskell}

Granule is a research language serving as a test bed for ideas
in graded types~\cite{DBLP:journals/pacmpl/OrchardLE19}. In our classification,
Granule is a linear-base system by default.

Semiring-graded modalities written here as $ \textcolor{coeffectColor}{\square_{ \SYSTEMnt{r} } }  \SYSTEMnt{A} $ in the
core calculi are written as `\granin{a [r]}' in Granule.
While Granule uses \granin{->}  for the function arrow (\texttt{->} in ASCII), as in Haskell or ML, it is really a linear arrow.
Our running example looks as follows in Granule:
%
\begin{granule}
const : forall {a : Type, b : Type, s : Semiring} . a -> b [0 : s] -> a
const x y = let [_] = y in x
\end{granule}
%
Granule also includes a language extension that changes its semantics to Graded
Base and whose syntax for graded arrows is borrowed from Linear Haskell, but supporting
semiring polymorphism:
%
\begin{granule}
language GradedBase
const' : forall {a : Type, b : Type, s : Semiring} . a -> b % (0 : s) -> a
const' x y = x
\end{granule}
%
%The graded-base approach has been explored as a target
%for program synthesis~\cite{DBLP:conf/esop/HughesO24}.


\subsection{Comparison with Effects}
\label{sec:comparison-with-effects}
%
Here we have considered the relationship between Linear Base
with linear function spaces and linear assumptions as well
as graded, and Graded Base where grades are pervasive. This
situation is somewhat analogous to the relationship between
\emph{effect systems} and \emph{monads}.

The type-and-effect systems of Gifford and
Lucassen~[\citeyear{gifford1986integrating,lucassen1988polymorphic}]
provide a technique for augmenting a type system with information
about the side effects a program may perform whilst computing
a result. Judgments and function arrows carry effect information:
%
\begin{align*}
\Gamma \vdash M : \tau, F \qquad\quad
\Gamma' \vdash \lambda x . M' : \sigma \xrightarrow{F'} \tau, \emptyset
\end{align*}
%
i.e., the program $M$ computes a result of type $\tau$ and may perform side effects described by
$F$. Traditionally $F$ is a set of effect `labels', describing
operations that have been performed.  This was later generalised to
semilattices and ordered
monoids~\cite{nielson1999type,DBLP:conf/birthday/MycroftOP16},
enabling other kinds of analysis. On the right,
the function term encapsulates side effects $F'$ and is itself pure,
marked with $\emptyset$. This is analogous to Graded Base, where coeffect grades annotate function types
(instead representing usage of the parameter).

Monads on the other hand provide a semantic technique for giving a
denotational model of various kinds of
effect~\cite{DBLP:journals/iandc/Moggi91,DBLP:conf/lics/Moggi89}.
The \emph{monadic meta language} provides a type constructor $T$
where $T \tau$ captures the type of computations performing some side
effects in order to compute a $\tau$ value~\cite{DBLP:journals/iandc/Moggi91}.
Thus, terms of the monadic meta language have types
of the form:
%
\begin{align*}
\Gamma \vdash M : T \tau \qquad\quad
\Gamma \vdash \lambda x . M' : \sigma \rightarrow T \tau
\end{align*}
%
On the right is an example of a function that encapsulates
side effects in its result, but the function itself does not need
wrapping in the monadic $T$ constructor since abstraction is pure
(side-effect free). This is analogous to Linear Base if we had
only a single comonadic modality.

\citet{wadler2003marriage} observed that the structure of effect systems
and the structure of the monadic meta language are the same.
They thus married the two by annotating the monadic constructor
with effect information, i.e.,
%
\begin{align*}
\Gamma \vdash M : T_F \tau \qquad\quad
\Gamma' \vdash \lambda x . M' : T_\emptyset (\sigma \rightarrow T_{F'} \tau)
\end{align*}
%
Later, this indexed structure was captured generally by the notion of
\emph{graded monads}~\cite{DBLP:conf/popl/Katsumata14,DBLP:journals/corr/OrchardPM14}. Graded
monads generalise monads to a family of constructors indexed by the
elements of a monoid, and whose operations are stratified and
accounted for by the monoidal structure. Graded monads thus unify
the semantics of effects with a syntactic analysis in the types.

Linear Base with graded modality $ \textcolor{coeffectColor}{\square_{ \SYSTEMnt{r} } }  \SYSTEMnt{A} $ is thus a dual
notion to the graded monadic calculus, unifying the semantics of
coeffects with a static analysis in the types, but with linearity
as its base typing system. We could thus see Linear Base as a \emph{graded comonadic meta language}.
In contrast, Graded Base is analogous to effect systems. Indeed the early
work on \emph{coeffect systems} (captured here as Graded Base) has no graded modality and was motivated
by dualising effect systems~\cite{DBLP:conf/icalp/PetricekOM13,petricek2014coeffects}.
As seen, without this graded modality, Graded Base cannot localise information
to sub-parts of complex types as discussed in Section~\ref{sec:prods-sums}.

\subsection{Adjoint Relationship}

\citet{DBLP:journals/corr/abs-2401-17199} follow
the Mixed Linear and Non-Linear logic approach of \citet{DBLP:conf/csl/Benton94}
to provide a pair of graded calculi with an adjunction mediating between them.
One calculus is a linear-base system with linear functions
and a mixed context of linear and graded assumptions. The other calculus
is a graded-base system but with no graded function arrow. A modality
$\mathsf{Lin}$ embeds formulae from the linear-base system into the graded-base. Conversely, a graded modality
$\mathsf{Grd}_r$ embeds formulae from the graded-base system into
the linear-base system, capturing their grade $r$ and incurring a promotion (scaling of
the graded context). From $\mathsf{Lin}$ and
$\mathsf{Grd}_r$ a graded comonadic modality is derived following the adjoint resolution of graded
comonads~\cite{DBLP:conf/fossacs/FujiiKM16}. %They then study the categorical semantics for this system.

The work seems related but on a more semantic level. Whilst we showed
that our translation is not adjoint in Section~\ref{sec:transl}, our two calculi have
function types on both sides which seemed to be the main blocker to forming an adjunction.
Further work is to ascertain whether our translation could
be restricted to the systems considered by \citet{DBLP:journals/corr/abs-2401-17199} and then used
as a model of the adjunction by using our
translations as the semantics of the modalities operating on
syntactic categories of terms. Another avenue
would be to explore whether generalising function spaces on the graded-base
side would undermine their adjunction construction in general.

\subsection{Filling Out the Design Space: Cartesian Base}
\label{sec:cartesian-base}

A central idea in this work has been to consider the `base' notion of usage in
graded systems: linear-base systems have linear propositions, with a linear
axiom rule, and linear function type; graded-base systems have graded propositions,
with a graded axiom rule, and a graded function type. This framing suggest
an unexplored point in the design space: what about a graded system with
a cartesian base, i.e., with a cartesian axiom rule and standard cartesian
function (type), with contexts comprising both cartesian and graded
assumptions?

Crucially we must then consider how the cartesian part interacts with
the graded part. On the left, we have the rules for dereliction and
application in Linear Base, on the right have the potential rules
for a Cartesian Base. Since cartesian propositions can be used arbitrarily
we now require an element $\omega \in \mathcal{R}$ of the semiring which
behaves in an absorbing fashion for both addition $+$ and multiplication $\cdot$.
Dereliction in Cartesian Base can then mark a cartesian assumption
with $\omega$. Application must also scale the argument context $\Gamma_2$
by $\omega$ since the function arrow is cartesian and so may use
its argument arbitrarily:

\begin{minipage}{0.5\linewidth}
\begin{align*}
\begin{array}{c}
\text{(Linear Base)} \\[0.5em]
\dfrac{\Gamma, x : A \vdash_{\textsc{l}} t : B}
      {\Gamma, x : [A]_1 \vdash_{\textsc{l}} t : B} (der)
\\[1.5em]
\dfrac{\Gamma_1 \vdash_{\textsc{l}} t_1 : A \multimap B \quad \Gamma_2 \vdash_{\textsc{l}} t_2 : A}
      {\Gamma_1 + \Gamma_2 \vdash_{\textsc{l}} t_1\ t_2 : B}(app)
\end{array}
\\
\end{align*}
\end{minipage}
\begin{minipage}{0.5\linewidth}
\begin{align*}
\begin{array}{c}
\text{(Cartesian Base)} \\[0.5em]
\dfrac{\Gamma, x : A \vdash t : B}
      {\Gamma, x : [A]_\omega \vdash t : B}(der)
\\[1.5em]
\dfrac{\Gamma_1 \vdash t_1 : A \rightarrow B \quad \Gamma_2 \vdash t_2 : A}
      {\Gamma_1 + \omega \ast \Gamma_2 \vdash t_1\ t_2 : B}(app)
\end{array}
\\
\end{align*}
\end{minipage}

\noindent
The above seems a reasonable approach to building a
Cartesian Base system. However, this system
provides little benefit as the grading quickly degenerates to
being always $\omega$ in derivations. Consider below the derivation of
the `copy' operation on an axiom to produce a pair:

\begin{align*}
\begin{array}{cc}
\dfrac{
\dfrac{x : A \vdash_{\textsc{l}} x : A}
      {x : [A]_1 \vdash_{\textsc{l}} x : A}
\quad
\dfrac{x : A \vdash_{\textsc{l}} x : A}
      {x : [A]_1 \vdash_{\textsc{l}} x : A}
}
{x : [A]_{1 + 1} \vdash_{\textsc{l}} (x, x) : A \otimes A}
&
\dfrac{
\dfrac{x : A \vdash x : A}
      {x : [A]_\omega \vdash x : A}
\quad
\dfrac{x : A \vdash x : A}
      {x : [A]_\omega \vdash x : A}
}
{x : [A]_\omega \vdash (x, x) : A \otimes A}
\end{array}
\end{align*}
%
In Linear Base, there is a precise account of the non-linearity
incurred: the term uses $x$ with graded $1+1$. However, in
Cartesian Base we can only say $\omega$ due to the dereliction rule.
In the end, the Cartesian Base system
is ultimately no more informative as a graded system than
standard intuitionistic, cartesian systems.


\subsection{Resource-Aware Operational Semantics}
\label{sec:raos}

The operational semantics presented here leverage standard $\lambda$-calculus CBN operational semantics,
which are typical for the literature studied.
More recent work has developed \emph{resource-aware} operational
semantics~\cite{grad, DBLP:journals/pacmpl/BianchiniDGZ23, DBLP:journals/corr/abs-2311-11795, DBLP:conf/ecoop/BianchiniDGZ23}. These approaches are based on an abstract machine
approach with a variable environment, rather than a substitution-based model.
Conceptually, a resource-aware semantics keeps track not just of variable mappings in the evaluation environment,
but more broadly tracks resource supplies, which are consumed by variable use,
thereby making the connection between static and dynamic resource guarantees:
A \emph{resource-sound} type system only accepts \emph{well-resourced} programs,
whose evaluation under the resource-aware operational semantics respects the static resource bounds,
meaning (1) evaluation does not get stuck due to depletion (\emph{resource progress});
and (2) every evaluation step leads to a program with comparable resource requirements
under the (transitive) approximation relation ($\sqsubseteq$) of the resource algebra (\emph{resource conservation}).

Most recently, \citet{DBLP:journals/corr/abs-2507-13792} extend these guarantees to the converse property,
ensuring both upper and lower resource bounds, thereby also preventing unused resources (\emph{wastefulness}).

Extending the various flavours of linear- and graded-base calculi and their translations,
which form the substrate of this paper, to such resource-aware operational semantics goes
beyond the scope of what we are able to investigate here. Instead, we leave this as future work.

\subsection{Conclusion}
\label{sec:conclusion}

Our goal here has been to serve as a bridge between two different approaches
in the literature for constructing graded type systems that capture coeffects (context dependence).
We answer the long open question about the relationship between
the linear-base and graded-base styles and give a detailed analysis and metatheoretic results. Overall we see that we can translate between
the two flavours: with Linear Base and Graded Base having mutual embeddings
(but requiring a global CPS encoding for the Linear Base into Graded Base direction);
Linear Base and Graded Modal Base having (more straightforward) mutual embeddings;
and Linear Push Core and Graded Modal Core having mutual embeddings: these are the points
where the two styles of calculi can be inter-translated (see Figure~\ref{fig:relationship}).

Now that graded types for tracking various kinds of program property are moving
from the laboratory to the field, with Haskell~\cite{DBLP:journals/pacmpl/BernardyBNJS18}, Idris~\cite{DBLP:conf/ecoop/Brady21}, Granule~\cite{DBLP:journals/pacmpl/OrchardLE19}, and lately OCaml with its moding system~\cite{DBLP:journals/pacmpl/LorenzenWDEL24,DBLP:journals/pacmpl/GeorgesPEWDECPD25},
our hope is that this study can serve as a design guide for
further theoretical work and for language designers wishing to integrate graded types
in some way, drawing from a wider arsenal of techniques and metatheoretic results.
